\answerpaper{数学 A 类}

\answerproblem{一、}{15}
设 $N(0,0,1)$ 是球面 $S:x^2+y^2+z^2=1$ 的北极点．$A(a_1,a_2,0)$，$B(b_1,b_2,0)$，$C(c_1,c_2,0)$ 为 $xOy$ 平面上不同的三点．设连接 $N$ 与 $A,B,C$ 的三直线依次交球面 $S$ 于点 $A_1,B_1$ 与 $C_1$．
\begin{enumerate}
  \item 求连接 $N$ 与 $A$ 两点的直线方程．
  \item 求点 $A_1,B_1$ 与 $C_1$ 的坐标．
  \item 给定点 $A(1,-1,0)$，$B(-1,1,0)$，$C(1,1,0)$，求四面体 $NA_1B_1C_1$ 的体积．
\end{enumerate}
\begin{solution}
\textbf{(1)} 过 $N,A$ 两点的直线方程为
\[
  \frac{x}{a_1}=\frac{y}{a_2}=\frac{z-1}{-1}.
\]

\textbf{(2)} 由此可得直线的参数方程
\[
  x=a_1t,\qquad y=a_2t,\qquad z=1-t,
\]
代入球面方程可得
\[
  (a_1t)^2+(a_2t)^2+(1-t)^2=1.
\]
由此解得
\[
  t=\frac{2}{a_1^2+a_2^2+1}\quad\text{或}\quad t=0.
\]
从而得 $A_1$ 的坐标为
\[
  \Bigl(\frac{2a_1}{a_1^2+a_2^2+1},\ \frac{2a_2}{a_1^2+a_2^2+1},\ \frac{a_1^2+a_2^2-1}{a_1^2+a_2^2+1}\Bigr).
\]
同理可得，$B_1$ 的坐标为
\[
  \Bigl(\frac{2b_1}{b_1^2+b_2^2+1},\ \frac{2b_2}{b_1^2+b_2^2+1},\ \frac{b_1^2+b_2^2-1}{b_1^2+b_2^2+1}\Bigr),
\]
以及 $C_1$ 的坐标为
\[
  \Bigl(\frac{2c_1}{c_1^2+c_2^2+1},\ \frac{2c_2}{c_1^2+c_2^2+1},\ \frac{c_1^2+c_2^2-1}{c_1^2+c_2^2+1}\Bigr).
\]

\textbf{(3)} 当 $A(1,-1,0)$，$B(-1,1,0)$ 以及 $C(1,1,0)$ 给定时，经计算可得
\[
  A_1=\Bigl(\frac{2}{3},-\frac{2}{3},\frac{1}{3}\Bigr),\qquad
  B_1=\Bigl(-\frac{2}{3},\frac{2}{3},\frac{1}{3}\Bigr),\qquad
  C_1=\Bigl(\frac{2}{3},\frac{2}{3},\frac{1}{3}\Bigr).
\]
所以，利用向量的混合积可以把四面体 $NA_1B_1C_1$ 的体积表示为
\[
  V=\frac{1}{6}\left|\bigl(\overrightarrow{NA_1},\overrightarrow{NB_1},\overrightarrow{NC_1}\bigr)\right|.
\]
混合积 $\bigl(\overrightarrow{NA_1},\overrightarrow{NB_1},\overrightarrow{NC_1}\bigr)$ 表示成矩阵的行列式
\[
  \bigl(\overrightarrow{NA_1},\overrightarrow{NB_1},\overrightarrow{NC_1}\bigr)
  =\det\begin{pmatrix}
    \frac{2}{3} & -\frac{2}{3} & \frac{2}{3}\\[2pt]
    -\frac{2}{3} & \frac{2}{3} & \frac{2}{3}\\[2pt]
    -\frac{2}{3} & -\frac{2}{3} & -\frac{2}{3}
  \end{pmatrix}
  =\frac{32}{27}.
\]
于是得到
\[
  V=\frac{1}{6}\times\frac{32}{27}=\frac{16}{81}.
\]
\end{solution}

\answerproblem{二、}{15}
求极限
\[
  \lim_{n\to\infty}\frac{\ln n}{\ln\bigl(1^{2020}+2^{2020}+\cdots+n^{2020}\bigr)}.
\]
\begin{solution}
我们有
\begin{align*}
&\lim_{n\to\infty}\frac{1}{n^{2021}}\bigl(1^{2020}+2^{2020}+\cdots+n^{2020}\bigr)\\
&\quad=\lim_{n\to\infty}\frac{1}{n}\Bigl(\Bigl(\frac{1}{n}\Bigr)^{2020}+\Bigl(\frac{2}{n}\Bigr)^{2020}+\cdots+\Bigl(\frac{n}{n}\Bigr)^{2020}\Bigr)\\
&\quad=\int_0^1 x^{2020}\dd x=\frac{1}{2021}.
\end{align*}
或用 Stolz 公式
\[
  \lim_{n\to\infty}\frac{1}{n^{2021}}\bigl(1^{2020}+2^{2020}+\cdots+n^{2020}\bigr)
  =\lim_{n\to\infty}\frac{n^{2020}}{n^{2021}-(n-1)^{2021}}=\frac{1}{2021}.
\]
因此，$\ln\dfrac{1^{2020}+2^{2020}+\cdots+n^{2020}}{n^{2021}}$ 有界．从而
\begin{align*}
&\lim_{n\to\infty}\frac{\ln n}{\ln(1^{2020}+2^{2020}+\cdots+n^{2020})}\\
&\quad=\lim_{n\to\infty}\frac{\ln n}{2021\ln n+\ln\frac{1^{2020}+2^{2020}+\cdots+n^{2020}}{n^{2021}}}\\
&\quad=\frac{1}{2021}.
\end{align*}
\end{solution}

\answerproblem{三、}{15}
设 $A,B$ 均为 2020 阶正交矩阵，齐次线性方程组 $Ax=Bx\ (x\in\R^{2020})$ 的解空间维数为 $3$．问：矩阵 $A,B$ 是否可能相似？证明你的结论．
\begin{solution}
$A,B$ 一定不相似．

证明如下．

令 $C=AB^{-1}$．由于 $A,B$ 均为正交矩阵，故 $C$ 也是正交矩阵．$C$ 视为复矩阵是酉矩阵，故可以复对角化．即存在复可逆矩阵 $T$ 和复对角矩阵 $D$，使得 $T^{-1}CT=D$，其中 $D$ 的主对角线上的元素即为 $C$ 的复特征值．

齐次线性方程组 $Ax=Bx$ 的解空间维数为 $3$，则 $\rank(A-B)=2017$．进而
\begin{align*}
\rank(D-I)&=\rank\bigl(T^{-1}(C-I)T\bigr)=\rank(C-I)\\
&=\rank\bigl((A-B)B^{-1}\bigr)=\rank(A-B)=2017.
\end{align*}
这表明对角矩阵 $D$ 的主对角线上恰有 $3$ 个元素是 $1$．即 $C$ 有三重特征值 $1$．由于正交矩阵的实特征值为 $1$ 或 $-1$，而非实数特征值共轭成对出现，共有偶数个．又 $C$ 共有 $2020$ 个特征值（计重数），故 $C$ 有特征值 $-1$，且重数为奇数．

$C$ 的行列式是其所有特征值（计重数）之积．注意到 $C$ 的非实数特征值共轭成对出现，它们的乘积为正数．故 $\det C<0$．特别地，$\det(AB^{-1})=\det C\neq 1$．即 $\det A\neq\det B$．从而 $A,B$ 不相似．
\end{solution}

\answerproblem{四、}{20}
称非常值一元 $n$ 次多项式（合并同类项后）的 $n-1$ 次项（可能为 $0$）为第二项．求所有 2020 次复系数首一多项式 $f(x)$，满足对 $f(x)$ 的每个复根 $x_k$，都存在非常值复系数首一多项式 $g_k(x)$ 和 $h_k(x)$，使得 $f(x)=(x-x_k)g_k(x)h_k(x)$，且 $g_k(x)$ 与 $h_k(x)$ 的第二项系数相等．
\begin{solution}
显然 $f(x)=x^{2020}$ 满足题意．

以下证明这是唯一解．设 $f(x)$ 的 $2020$ 个复根为 $x_1,x_2,\cdots,x_{2020}$．对每个 $k\ (1\leqslant k\leqslant 2020)$，由题设条件可设 $f(x)=(x-x_k)g_k(x)h_k(x)$，其中 $g_k(x),h_k(x)$ 分别为 $m_k$ 次和 $n_k$ 次非常值首一多项式，第二项系数均为 $a_k$．设 $g_k(x)$ 的所有复根为 $y_{k,1},y_{k,2},\cdots,y_{k,m_k}$，$h_k(x)$ 的所有复根为 $z_{k,1},z_{k,2},\cdots,z_{k,n_k}$．这些根恰为所有 $x_j\ (j\neq k)$．由韦达定理，
\[
  a_k=(-1)^{m_k-1}\bigl(y_{k,1}+y_{k,2}+\cdots+y_{k,m_k}\bigr)
  =(-1)^{n_k-1}\bigl(z_{k,1}+z_{k,2}+\cdots+z_{k,n_k}\bigr).
\]

对每个 $k$，将上式改写为
\[
  \sum_{j\neq k}\varepsilon_{kj}x_j=0,
\]
其中 $\varepsilon_{kj}=1$ 或 $-1$．

这样，我们得到了关于 $x_1,x_2,\cdots,x_{2020}$ 的齐次线性方程组，其系数矩阵 $A$ 为 $2020$ 阶方阵，主对角线上元素为 $0$．主对角线外元素为 $1$ 或 $-1$．令 $B$ 为 $2020$ 阶方阵，其主对角线上元素为 $0$，主对角线外元素为 $1$，则容易计算出 $B$ 的行列式为 $\det B=-2019$．由行列式定义，$\det A$ 与 $\det B$ 的奇偶性相同，故 $\det A\neq 0$．

从而上述齐次线性方程组只有零解，即 $x_1=x_2=\cdots=x_{2020}=0$．这便证明了 $f(x)=x^{2020}$．

\textbf{注：}上面证明 $\det A\neq 0$ 也可以如下进行：

显然，$\det A\equiv\det B\pmod 2$．由于 $\det B=-2019\equiv 1\pmod 2$．所以 $\det A\equiv 1\pmod 2$．故 $\det A\neq 0$．
\end{solution}

\answerproblem{五、}{15}
设 $\varphi$ 是 $\R$ 上严格单调增加的连续函数，$\psi$ 是 $\varphi$ 的反函数，实数列 $\{x_n\}$ 满足
\[
  x_{n+2}=\psi\Bigl(\Bigl(1-\frac{1}{\sqrt{n}}\Bigr)\varphi(x_n)+\frac{1}{\sqrt{n}}\varphi(x_{n+1})\Bigr),\qquad n\geqslant 2.
\]
证明 $\{x_n\}$ 收敛或举例说明 $\{x_n\}$ 有可能发散．
\begin{solution}
我们断言 $\{x_n\}$ 收敛．证明如下．记 $y_n=\varphi(x_n)$，则
\[
  y_{n+2}=\Bigl(1-\frac{1}{\sqrt{n}}\Bigr)y_n+\frac{1}{\sqrt{n}}y_{n+1},\qquad n\geqslant 2.
\]
令
\[
  a_n=\min\{y_n,y_{n-1}\},\qquad b_n=\max\{y_n,y_{n-1}\},\qquad n\geqslant 3.
\]
则
\[
  a_n\leqslant y_{n+1}\leqslant b_n,\qquad n\geqslant 3.
\]
进而
\[
  a_n\leqslant a_{n+1}\leqslant b_{n+1}\leqslant b_n,\qquad n\geqslant 3.
\]
所以 $\{a_n\},\{b_n\}$ 均单调有界，从而收敛．

特别，$\{y_n\}$ 有界．由于
\[
  y_{n+2}-y_{n+1}=-\Bigl(1-\frac{1}{\sqrt{n}}\Bigr)(y_{n+1}-y_n),\qquad n\geqslant 2.
\]
因此
\[
  |y_{n+2}-y_{n+1}|\leqslant |y_3-y_2|\prod_{k=2}^{n}\Bigl(1-\frac{1}{\sqrt{k}}\Bigr),\qquad n\geqslant 2.
\]
由 $\prod\limits_{k=2}^{\infty}\Bigl(1-\dfrac{1}{\sqrt{k}}\Bigr)$ 发散到零，得到 $\lim\limits_{n\to\infty}(y_{n+1}-y_n)=0$．

所以
\[
  b_n-a_n=|y_n-y_{n-1}|\to 0,\qquad n\to\infty.
\]
这样，$\{a_n\}$ 与 $\{b_n\}$ 的极限相等，从而 $\{y_n\}$ 收敛．

最后，由 $\psi$ 的连续性得到 $\{x_n\}$ 收敛．
\end{solution}

\answerproblem{六、}{20}
对于有界区间 $[a,b]$ 的划分 $P:a=x_0<x_1<\cdots<x_{n+1}=b$，其范数定义为 $\|P\|=\max\limits_{0\leqslant k\leqslant n}(x_{k+1}-x_k)$．现设 $[a,b]$ 上函数 $f$ 满足 Lipschitz 条件，即存在常数 $M>0$ 使得对任何 $x,y\in[a,b]$，成立 $|f(x)-f(y)|\leqslant M|x-y|$．定义
\[
  s(f;P)\equiv\sum_{k=0}^{n}\sqrt{|x_{k+1}-x_k|^2+|f(x_{k+1})-f(x_k)|^2}.
\]
若 $\lim\limits_{\|P\|\to 0^+}s(f;P)$ 存在，则称曲线 $y=f(x)$ 可求长．记 $P_n$ 为 $[a,b]$ 的 $2^n$ 等分．证明：
\begin{enumerate}
  \item $\lim\limits_{n\to\infty}s(f;P_n)$ 存在．
  \item 曲线 $y=f(x)$ 可求长．
\end{enumerate}
\begin{solution}
\textbf{法 I.} 我们有
\[
  0\leqslant s(f;P)\leqslant\sum_{k=0}^{n}\sqrt{M^2+1}\,|x_{k+1}-x_k|=(b-a)\sqrt{M^2+1}.
\]
因此，$s(f;P)$ 有界．

\textbf{(1)} 由平面上点和点距离的三角不等式，立即有
\[
  s(f;P_n)\leqslant s(f;P_{n+1}),\qquad \forall\, n\geqslant 1.
\]
因此，$\{s(f;P_n)\}$ 单调增加，结合有界性知其收敛．设极限为 $L$．

一般地，对于划分 $P,Q$，用 $P\oplus Q$ 表示由 $P$ 和 $Q$ 的所有分点为分点的划分，则
\[
  s(f;P\oplus Q)\geqslant s(f;P).
\]

\textbf{(2)} 对于任何 $\varepsilon>0$，有 $m\geqslant 1$ 使得
\[
  s(f;P_m)\geqslant L-\varepsilon.
\]
对于划分 $P$，用 $P\oplus P_m$ 表示由 $P$ 和 $P_m$ 的所有分点为分点的划分，则
\[
  s(f;P\oplus P_m)\geqslant s(f;P_m)\geqslant L-\varepsilon.
\]
在 $s(f;P\oplus P_m)$ 的和式中，与 $s(f;P)$ 的和式中不同的项是涉及 $P_m$ 的分点的项，总数不超过 $2^{m+1}$ 项，相应的小区间长度不超过 $\|P\oplus P_m\|\leqslant\|P\|$．因此，这些项的和不超过 $2^{m+1}\sqrt{M^2+1}\,\|P\|$．于是
\begin{align*}
s(f;P)&\geqslant s(f;P\oplus P_m)-2^{m+1}\sqrt{M^2+1}\,\|P\|\\
&\geqslant L-\varepsilon-2^{m+1}\sqrt{M^2+1}\,\|P\|.
\end{align*}
这样
\[
  \varliminf_{\|P\|\to 0^+}s(f;P)\geqslant L-\varepsilon.
\]
进而
\[
  \varliminf_{\|P\|\to 0^+}s(f;P)\geqslant L.
\]

类似地，记 $K=\varlimsup\limits_{\|P\|\to 0^+}s(f;P)$．对于任何 $\varepsilon>0$，有划分 $Q$ 使得
\[
  s(f;Q)\geqslant K-\varepsilon.
\]
则
\[
  s(f;Q\oplus P_m)\geqslant s(f;Q)\geqslant K-\varepsilon.
\]
在 $s(f;Q\oplus P_m)$ 的和式中，与 $s(f;P_m)$ 的和式中不同的项是涉及 $Q$ 的分点的项，总数不超过 $2N$ 项，其中 $N$ 是划分 $Q$ 的分点数．因此，这些项的和不超过 $2N\sqrt{M^2+1}\,\|P_m\|$．于是
\begin{align*}
s(f;P_m)&\geqslant s(f;Q\oplus P_m)-2N\sqrt{M^2+1}\,\|P_m\|\\
&\geqslant K-\varepsilon-2N\sqrt{M^2+1}\,\|P_m\|.
\end{align*}
这样
\[
  L=\lim_{m\to\infty}s(f;P_m)\geqslant K-\varepsilon.
\]
进而 $L\geqslant K$．结合 $K\geqslant\varliminf\limits_{\|P\|\to 0^+}s(f;P)\geqslant L$ 得到
\[
  \lim_{\|P\|\to 0^+}s(f;P)=L.
\]
即 $y=f(x)$ 可求长．

\textbf{法 II.} 事实上，注意到 (1) 是 (2) 的推论，我们只需直接证明 (2)．具体证明如下．

我们有
\[
  0\leqslant s(f;P)\leqslant\sum_{k=0}^{n}\sqrt{M^2+1}\,|x_{k+1}-x_k|=(b-a)\sqrt{M^2+1}.
\]
因此，$s(f;P)$ 有界．

对于划分 $P,Q$，用 $P\oplus Q$ 表示由 $P$ 和 $Q$ 的所有分点为分点的划分，由平面上点和点距离的三角不等式，立即有
\[
  s(f;P\oplus Q)\geqslant s(f;P).
\]

考虑划分列 $\{Q_k\}$ 使得 $\lim\limits_{k\to\infty}\|Q_k\|=0$，且
\[
  \lim_{k\to\infty}s(f;Q_k)=L=\varlimsup_{\|P\|\to 0}s(f;P).
\]

对于每个 $k\geqslant 1$，设 $N_k$ 为划分 $Q_k$ 的分点数．
\[
  s(f;P\oplus P_m)\geqslant s(f;P_m)\geqslant L-\varepsilon.
\]
{\setlength{\emergencystretch}{1em}%
在 $s(f;P\oplus Q_k)$ 的和式中，与 $s(f;P)$ 的和式中不同的项是涉及 $Q_k$ 的分点的项，总数不超过 $2N_k$ 项，相应的小区间长度不超过 $\|P\oplus Q_k\|\leqslant\|P\|$．因此，这些项的和不超过 $2N_k\sqrt{M^2+1}\,\|P\|$．于是\par}
\begin{align*}
s(f;P)&\geqslant s(f;P\oplus Q_k)-2N_k\sqrt{M^2+1}\,\|P\|\\
&\geqslant s(f;Q_k)-2N_k\sqrt{M^2+1}\,\|P\|.
\end{align*}
这样
\[
  \varliminf_{\|P\|\to 0^+}s(f;P)\geqslant s(f;Q_k),\qquad \forall\, k\geqslant 1.
\]
进而
\[
  \varliminf_{\|P\|\to 0^+}s(f;P)\geqslant L=\varlimsup_{\|P\|\to 0^+}s(f;P).
\]
所以 $\lim\limits_{\|P\|\to 0^+}s(f;P)$ 存在．即 $y=f(x)$ 可求长．自然也有 $\lim\limits_{n\to\infty}s(f;P_n)$ 存在．
\end{solution}
